\documentclass[11pt,a4paper]{article}

% ===== Packages =====
\usepackage[utf8]{inputenc}
\usepackage[T1]{fontenc}
\usepackage{geometry}
\geometry{margin=1in}
\usepackage{hyperref}
\usepackage{enumitem}
\usepackage{microtype}

% This extended abstract carries no equations. It states the framework as a
% research program and leaves each mathematical construction open, for new
% collaborators. Information geometry is retained as the one standard foundation.
%
% Authorship to be confirmed.

\title{Alignment as the Geometry of Belief-Generating Systems\\
\large An Extended Abstract and Research Program}

\author{[authors to be confirmed]}
\date{}

\begin{document}

\maketitle

\begin{abstract}
Most current measures of AI alignment evaluate what an agent produces rather
than how it produces it. Reinforcement learning from human feedback evaluates
behavior; information-theoretic methods evaluate belief distributions. Neither
captures the dynamics that generate the beliefs in the first place. Two agents
can hold the same belief now and still respond differently to new evidence,
social pressure, or manipulation, because one holds it rigidly and the other
loosely. We propose treating alignment as a relationship between the systems
that generate beliefs, described by their attractor landscapes, the paths
beliefs follow through those landscapes, and the geometry of the landscape
itself. This abstract sets out the framework, the phenomena it should account
for, and the patterns we have already found in real data. We keep one standard
mathematical foundation, information geometry, which gives the space of beliefs
a curved geometry. The other objects we need, a structural distance between
belief dynamics, a least-effort belief transition, and a measure of landscape
curvature, are left as open problems rather than fixed definitions.
\end{abstract}

\section{The problem}

Current alignment methods work almost entirely on what an agent outputs.
Preference-based training adjusts behavior to match human choices,
constitution-based methods impose principles stated as text, and quantitative
work usually measures the divergence between an agent's output distribution and
a target.

What these methods miss is the process behind the output. Suppose two agents
currently assign the same probability to a proposition. Output-level and
distributional measures call them equally aligned. But one agent may hold the
belief firmly, with a conviction that resists revision and returns after a
perturbation, while the other holds it loosely and updates readily. Given new
information, a deliberate nudge, or a change in their social setting, the two
agents behave very differently. The first is fragile: aligned now, possibly
badly misaligned later. The second can track a changing world. Existing
measures cannot separate the two cases, because the difference is in the
dynamics and not in the current belief.

\section{The thesis}

\begin{quote}
\emph{Alignment should be understood as the distance between belief-generating
dynamical systems, described by their attractor landscapes, the trajectories
beliefs trace through them, and the curvature of the terrain, rather than as the
distance between the beliefs those systems currently produce.}
\end{quote}

This changes several things at once. What we measure becomes the shape of the
belief-forming process instead of its momentary output. Good alignment becomes
agreement that survives perturbation. Intervention becomes the reshaping of the
landscape an agent moves through, which is a different act from overwriting the
conclusions it currently holds.

\section{The framework}

The program has five stages. None commits to a particular formalism here; each
identifies a phenomenon and the kind of mathematical object that would capture
it.

\paragraph{1. The landscape of belief.}
An agent that updates its beliefs against a stream of experience does not move
at random. Its beliefs tend to settle into recurring configurations, the
patterns its reasoning returns to. We call the set of these patterns the
agent's belief landscape and the stable patterns its attractors. The starting
point is information geometry, which is standard and which we keep. An agent's
beliefs are probability distributions, and the space of probability
distributions has a curved Riemannian geometry in which the distance between
two beliefs reflects how hard they are to tell apart statistically. This
supplies the substrate the landscape sits on: a curved manifold of belief
states with a canonical metric, so that distance, gradient, and curvature are
already defined. We take that as given. What concerns us is how an agent forms
and holds beliefs, not only what it currently believes, since the same stated
opinion can sit in landscapes of very different shape, and the shape is what
governs behavior under pressure. The first task is to recover the landscape
from observed behavior: its basins, its connectivity, and its complexity, using
trajectories alone and without access to the agent's internals.

\paragraph{2. What the landscape means.}
The geometry has a psychological reading. The depth of a basin reflects the
cognitive, emotional, and social effort needed to give up a belief; deep basins
are dogmas and shallow ones are working hypotheses. The ridges between basins
act as identity barriers, since crossing one can mean leaving a group, a
self-image, or a social world, so a ridge's height measures the real cost of
changing one's mind rather than the logical distance between positions. Funnels,
where many trajectories converge, correspond to shared narratives that draw
different agents toward a common view. This reading yields testable
predictions: hesitation and arousal near ridge crossings, resistance to belief
change that scales with basin depth, and social clustering that follows basin
membership.

\paragraph{3. Measuring misalignment structurally.}
This is the central problem and the one most in need of new mathematics. We
want a distance between two agents that reflects how differently their beliefs
evolve, not where their beliefs currently sit. Several requirements constrain
it. Two agents whose current beliefs match but whose landscapes have opposite
shape, one rigid and one flexible, should come out far apart. Two agents that
respond to evidence the same way should come out close, even when their
momentary opinions differ. The distance should be stable under noise and
short-lived fluctuations, so that it reflects lasting behavior. And it should
obey the usual axioms of a distance, so that the separations among three or
more agents stay mutually consistent. We state this as a requirement, not a
definition. Constructing such a distance between belief-generating systems is
the main problem we want to work on.

\paragraph{4. The paths beliefs take, and steering them.}
Beliefs change along trajectories through the landscape, and some trajectories
are more likely than others. There is a least-effort path: the most probable
route from one worldview to another, the one that works least against the
terrain. Knowing these paths is what makes steering possible. The aim is not to
force an agent across a ridge but to learn which interventions lower a ridge,
widen a valley, or deepen a basin, and how each changes the routes open to the
agent. Belief change also runs on several timescales at once, from fast
perceptual updates through slower revisions of opinion to slow shifts in core
identity, and the level that matters for comparing and steering agents is the
slow one, where the durable structure is. Two agents that differ only in fast
fluctuations but share a slow landscape are aligned in the sense we care about.
Formalizing the least-effort path and the separation of timescales are two
further open problems.

\paragraph{5. Which landscapes are desirable.}
Measuring and steering belief landscapes forces a normative question: which
landscapes should we steer toward? We suggest four properties of a good one. It
should be truth-tracking, with attractors that move under genuine evidence and
hold up against noise and manipulation. It should preserve flexibility, leaving
agents able to revise their beliefs instead of being trapped in a basin; this
is the geometric counterpart of corrigibility. It should be cooperatively
compatible, keeping bounded the cost of reaching mutual understanding among
agents in a coalition. And it should be pluralistic, supporting many distinct
but mutually reachable basins, so that agents can hold different worldviews and
still count as aligned as long as the routes between them stay affordable.
Alongside these we place constraints on intervention: lower ridges rather than
push agents over them, keep interventions reversible, share their costs fairly,
and act only at a level the agent can understand and consent to.

\section{Empirical support}

The quantities the framework posits can be measured, and we have found them in
three different real systems using the same analysis pipeline in each.

\begin{itemize}[leftmargin=1.4em]
\item \textbf{Rigid and flexible attractors can be told apart.} In a controlled
belief-network model with known ground truth, populations tuned to be rigid
versus flexible produced attractor landscapes whose topological complexity
differed by more than a factor of two, with the flexible populations forming
richer and more diffuse structure. Run on real financial market data, the same
pipeline still separated higher-volatility from lower-volatility regimes by the
shape of their dynamics.

\item \textbf{Cooperation and competition appear in the curvature of
interaction.} From dual-brain EEG recordings of eight pairs of people doing
cooperative and competitive tasks, a curvature measure on the moment-to-moment
coupling between their brain signals tracked which task each pair was
performing, in all eight pairs. The sign of the effect varied across pairs,
reflecting individual differences in how people couple, but the curvature
carried the cooperation/competition signal throughout.

\item \textbf{Echo chambers and open forums have different shape.} The same
topological analysis, applied to language embeddings from real online political
communities, found denser and more persistent structure in ideologically
homogeneous and polarized communities than in communities organized around
viewpoint diversity. The belief basins of echo chambers are measurably tighter.
\end{itemize}

The dynamics-based view repeatedly separates cases that a purely distributional
view merges, but it needs more data to do so. The structural comparison
requires long, rich behavioral trajectories, and where those are missing it
falls back toward distributional methods. Lowering that data requirement is one
of the practical goals for the new work.

\section{Open problems}

We close by listing the formal objects we need and do not yet have.

\begin{enumerate}[leftmargin=1.6em]
\item \textbf{A structural distance between belief-generating systems.} A
metric that compares two agents by how their beliefs evolve, stays stable under
noise, weighs lasting behavior over transient behavior, and can be computed
from observed trajectories. It can build on the information-geometric substrate
we keep, but it has to measure the dynamics on that manifold, not the distance
between two points on it. It should improve on any purely distributional
comparison.

\item \textbf{Recovering the landscape from behavior.} A stable method for
reconstructing attractors, basins, ridges, and connectivity from observed
trajectories, with guarantees that the result is robust to measurement noise.

\item \textbf{The least-effort belief transition.} A description of the most
probable path an agent takes between worldviews, the cost of that path, and how
interventions change it. This is the basis for steering that does not coerce.

\item \textbf{Timescale separation.} A rigorous account of how to average out
fast belief fluctuations and isolate the slow landscape at which alignment
should be measured and steered.

\item \textbf{From flexibility to a guarantee.} A formal connection between an
agent's capacity to revise its beliefs and a bound on how reliably it can be
steered to a target, so that corrigibility becomes something one can prove.

\item \textbf{Normative geometry.} Turning the four criteria for a good
landscape into precise, checkable conditions, and working out the trade-offs
among them.
\end{enumerate}

Our working assumption is that the right language for alignment is geometric:
basins, ridges, paths, curvature, and an agent's capacity to revise. Developing
that language carefully is what would let us state alignment precisely enough to
measure it and to prove things about it.

\end{document}
